\chapter{股票市场与股权请求权}
\label{chap:Stock_Markets_and_Equity_Claims}

\begin{openingcase}{为什么在没有现金交换的情况下，市场价值仍然会发生变化？}

一家上市公司有~1~亿股已发行在外股份，股价为人民币~10~元，股权市值为~10~亿元。某日早晨，公司宣布一项新产品取得重要技术进展。公司未支付股利、未发行新股，也没有收到现金，但投资者调整了对未来利润的预期，股价升至~12~元。

公司股权市值增加了人民币~2~亿元，却并未获得~2~亿元新增融资。这一变化发生在投资者相互转让现有股份的二级市场中。

\end{openingcase}

\begin{tcolorbox}[
  enhanced,
  colback=FinanceBlue!6,
  colframe=FinanceBlue,
  boxrule=0.6pt,
  sharp corners,
  left=8pt,
  right=8pt,
  top=7pt,
  bottom=7pt,
  before skip=1em,
  after skip=1.6em,
  title=\textbf{核心问题},
  coltitle=white,
  colbacktitle=FinanceBlue,
  fonttitle=\bfseries,
  attach boxed title to top left={xshift=6pt,yshift=-2mm},
  boxed title style={
    boxrule=0pt,
    sharp corners,
    left=5pt,
    right=5pt,
    top=2pt,
    bottom=2pt
  }
]
\begin{enumerate}[
  leftmargin=*,
  label=\textcolor{FinanceBlue}{\textbf{\arabic*.}},
  itemsep=0.55em,
  topsep=0.2em
]

\item 股权请求权提供哪些权利？剩余现金流权、有限责任以及控制权与债务所附带的权利有何不同？

\item 所有权结构、投票权以及公司治理机制如何影响着控制权分配，以及股东与管理者之间的代理冲突呢？

\item 在股票市场中，股权请求权是如何发行和交易的？信息、增长预期以及风险又如何反映在市场价格和必要报酬率中呢？

\end{enumerate}
\end{tcolorbox}

\section{作为剩余所有权请求权的股权}

股权是对公司的剩余请求权。股东并不持有获得固定支付的契约承诺，其权利价值取决于企业履行契约义务后的剩余价值，以及股份附带的治理权利。

\subsection{剩余现金流权与有限责任}

公司是独立法律实体。股份具有由公司法、公司章程及相应股份类别条款界定的财务与治理权利。普通股权对公司现金流和资产享有剩余请求权。清算时，普通股股东排在债权人和其他优先请求权人之后。

对于已缴足股款的股份，有限责任通常意味着股东不因持有股份而必须额外投入个人资产。但个人担保、借款、衍生品或其他契约产生的独立义务，仍可依法执行。

剩余地位并不意味着立即分配。盈利公司可以留存收益并再投资，而不向股东支付。股东对这些资源所产生的价值继续享有剩余请求权，其价值受未来经营结果及约束资源使用的治理安排影响。股利是向股东分配价值的一种方式，并非剩余请求权的定义。

\subsection{股权融资与风险承担}

股权融资提供资金，却不产生固定的契约利息或本金支付。公司发行普通股，以未来剩余现金流请求权及通常附带的投票权换取资本。经营现金流不确定时，这种支付灵活性可能具有价值。

不存在固定支付义务，同时意味着承担剩余风险。经营表现恶化时，普通股股利可能减少或取消，股权价值可能大幅下降；经营表现良好时，优先请求权得到满足后，股东分享上涨收益。因此，在考虑契约请求权后，普通股权承担企业价值变化的很大一部分风险。

这种承担风险的结构能够支持对不确定项目的投资。虽然仍然存在潜在的业务风险，但融资结构会影响负面经营结果如何传递至不同的利益相关者手中。

\begin{example}{商业风险与风险承担}
{business-risk-bearing}

假设一家公司正在开发一项新技术，初始投资为1亿元人民币，其中4000万元人民币用于债务融资，6000万元人民币用于股权融资。假设债务契约要求在债务到期时偿还4400万元人民币。假设项目所产生的收益足以满足这些请求权，且该公司没有其他资产或负债。

如果该项目在到期时产生1.4亿元人民币的收入，那么债务持有人将获得约定的4400万元人民币，而股东则能获得剩余的9600万元人民币。如果项目价值降至6000万元人民币，债务持有人仍然能收到4400万元人民币的回报，这样股东就能获得1600万元人民币的收益。然而，如果项目价值进一步降至3000万元人民币，股东将失去他们的全部投资，而债务持有人则只能收回3000万元人民币。


\end{example}

\subsection{债务与股权：优先次序、控制与风险分配}

债务和股权是对同一企业不同的请求权。它们的区别不仅仅体现在支付方式上。它们在优先权、控制权以及对企业价值变化的敏感度方面都存在差异。

债务赋予债权人契约规定的请求权，以及相对于普通股权的优先地位。普通债务的上行收益通常限于契约支付额，承诺支付未履行或未足额回收时则产生损失。债权人通常不行使一般公司控制权，但契约条件被违反或企业陷入财务困境时，契约限制条款、执行权利及重组安排可能赋予其更大影响力。



\begin{table}[htbp]
\centering

\caption{债务和股权作为金融请求权}
\label{tab:ch9-debt-equity-claims}

\begin{tabularx}{\textwidth}{@{}p{3.0cm}Y Y@{}}

\toprule

\textbf{维度}
& \textbf{债务}
& \textbf{普通股权} \\

\midrule

现金流权利
& 根据合同规定的利息和/或本金支付，在违约情况下执行
& 在满足合同义务后，对企业现金流的剩余请求权 \\

优先级&通常，清偿顺位高于普通股权；债务请求权的优先级取决于清偿顺位、证券以及适用法律的规定&清算中的剩余索赔人\\

到期日  
& 通常会有明确的指定日期，不过也存在永久性的债务安排。  
& 没有固定的偿还期限  
\\

控制项  
&通常缺乏普通的表决权；合同保护及执行权利遵循债务协议中的规定  
& 表决权与治理权，受股份类别及适用法律的约束  
\\

Upside  
& 通常仅限于合同金额的支付  
& 参与剩余收益的实现以及企业价值的提升  
\\

缺点
& 如果合同中的请求权未能完全得到偿还，可能会遭受损失
& 投资可能会全部损失，但股东通常仅因持有股份而无需对公司的债务承担个人责任 \\

\bottomrule

\end{tabularx}

\end{table}

普通股权是剩余请求权。股东没有约定的到期支付，在考虑优先请求权后获得剩余价值。在有限责任下，股份价值可能降至零，而上行收益不受契约偿还金额封顶。正常情况下，投票权还赋予普通股股东治理权利，包括选举董事和批准特定公司事项。

这些差异可能造成债权人与股东的利益冲突。债务存续期间，股东可能偏好上行收益更大的项目，即使额外下行损失部分由债权人承担。债权人通过定价、抵押品、契约限制条款、期限条款及其他契约保护作出应对。财务困境中，议价能力和控制权也可能向债权人转移。


\section{所有权、控制与公司治理}

公司形式将所有权与日常管理分离开来。股东拥有剩余请求权，而董事和经理则对公司的资源行使委托控制权。这种分离可能会在决策者无法承担其决策的全部后果时引发代理问题。公司治理机制负责分配控制权，并提供监督和惩戒管理的机制。

\subsection{投票权与董事会控制}

股东通常不会做出日常运营决策。他们主要通过投票权来行使控制权，包括选举董事和批准特定的重要事项。董事会任命并监督高级管理层，同时行使根据公司法以及公司治理安排所赋予的决策权。因此，董事会是股东投票权与管理层控制之间的重要纽带。

图~\ref{fig:ch9-governance-chain}展示了一个简化的委托控制链。

\begin{figure}[htbp]
\centering

\begin{tikzpicture}[
  >=Latex,
  box/.style={
    draw,
    rounded corners,
    align=center,
    minimum width=4.2cm,
    minimum height=0.85cm,
    inner xsep=6pt,
    inner ysep=4pt,
    font=\normalsize
  },
  flow/.style={
    -Latex,
    line width=0.85pt,
    FinanceBlue
  },
  lab/.style={
    font=\small,
    text=FinanceBlue,
    fill=white,
    inner xsep=3pt,
    inner ysep=1pt
  },
  node distance=0.75cm
]

\node[box] (share) {股东们};

\node[box, below=of share] (board)
  {董事会成员};

\node[box, below=of board] (mgmt)
  {管理};

\node[box, below=of mgmt] (assets)
  {企业资产与运营};

\draw[flow]
  (share) --
  node[lab, right] {选举董事}
  (board);

\draw[flow]
  (board) --
  node[lab, right] {指定并监督实施。}
  (mgmt);

\draw[flow]
  (mgmt) --
  node[lab, right] {保留LaTeX命令、公式、交叉引用和数字。}
  (assets);

\end{tikzpicture}

\caption{简化的委托控制链}
\label{fig:ch9-governance-chain}
\end{figure}

在一股一票结构下，正式投票权与持有的有投票权股份数量成正比。因此，持有~30\%~已发行在外有投票权股份的股东，也持有相应~30\%~的投票权。但实际影响力还取决于所有权集中度、投票参与程度、投票协议、董事选举规则及需股东批准的事项。

在表决权具有经济价值的情况下，股东能够影响公司的决策制定。即使某位股东的持股比例较低，只要公司股权分散，而其他股东参与投票的积极性不高，该股东仍然可以拥有相当大的影响力。因此，表决权的分配并不一定意味着实际控制权的分配也是如此。

\subsection{代理问题与所有权结构}

当决策权被下放时，就会出现代理问题，因为决策者与利益相关者的利益往往并不一致。在上市公司中，股权结构会影响对 corporate 决策的监督力度，同时也会引发由控制权带来的各种冲突。其中两种关系尤为重要。

第一种情况发生在管理者和股东之间。管理者掌控着公司的资源，但他们可能只拥有少量的剩余现金流权利。因此，他们的目标可能与股东的目标有所不同。管理者可能会重视私人利益、组织扩张、更高的工作稳定性，或者更低的个人风险，即使这些选择会降低股东的价值。

分散化的所有权结构会削弱监督激励机制的作用。小股东可能承担大部分对管理层进行调查的成本，而他们获得的公司价值提升却只占很小的一部分。其他股东则可以从中受益，而无需承担任何成本。这种“搭便车”现象会降低个人进行监督的积极性，同时也会增强董事会、机构投资者、信息披露、审计以及其他治理机制的重要性。

第二种关系涉及控制性股东与少数股东权益之间的关系。当股权高度集中时，需要加强对这种关系的监管。大型股东能够从中获得更多的利益，因此他们更有动力来监督管理层的行为。然而，过度集中的股权结构可能会引发控制性股东与少数股东之间的冲突。控制性股东可能会通过影响关联方交易、资产转让、融资决策以及其他企业行为等方式，谋取私利，同时降低少数投资者的价值。


\begin{table}[H]

\centering

\caption{两个重要的股权代理问题}

\label{tab:ch9-agency-problems}

\begin{tabularx}{\textwidth}{@{}p{3.0cm}Y Y@{}}

\toprule

\textbf{代理问题} & \textbf{冲突的来源} & \textbf{潜在的治理措施} \\

\midrule

管理者与股东&管理者拥有决策权，但只能持有部分剩余请求权&董事会监督、薪酬设计、信息披露、监控、投票权、收购压力\\

控制性股权与少数股东
& 控制性股权持有者能够影响决策，而这些决策的影响同样会被少数投资者所承担
& 关联方交易规则、独立董事制度、信息披露要求、投票保护机制以及法律保护 \\

\bottomrule

\end{tabularx}

\end{table}


\subsection{股东权利与股份类别}

一份股份代表着一系列财务和治理权利。特定股份所附带的权利取决于公司法律、公司的管理文件以及相关的股份类别。不同股份类别之间的差异可能会影响财务利益的分配以及投票权的分配。

不同股份类别可以不同方式分配这些权利。例如，优先股可能在股利或清算分配中优先于普通股，但投票权有限或没有投票权；其地位仍次于债权人请求权。转换、赎回或其他契约特征，可使优先股兼具股权和债务的某些特点。

存续的二元类和加权投票结构更直接地将投票权力与经济所有权分离开来。假设类\(j\)有\(N_j\)份存续股份，每股拥有\(v_j\)票。该类的总投票权比例为

\begin{equation}
\omega_j
=
\frac{v_jN_j}{\sum_k v_kN_k}.
\label{eq:ch9-voting-power}
\end{equation}

如果所有班级的每股现金流权益都相同，那么班级 \(j\) 所拥有的现金流份额为

\begin{equation}
\alpha_j
=
\frac{N_j}{\sum_k N_k}.
\label{eq:ch9-cash-flow-share}
\end{equation}

当不同类别的投票权存在差异时，\(\omega_j\)和\(\alpha_j\)不必相等。

这样的结构可以让创始人或其他控股股东在从外部投资者那里筹集股权的同时，仍然保持投票控制权。这可能会支持企业战略的连续性，尤其是在投资项目具有长期前景的情况下。同时，这种分离可能会削弱对外部股东的问责制，因为这些股东的经济影响力可能超过他们的投票影响力。治理方面的影响则取决于公司法律和上市规则所提供的保障措施。

\begin{tcolorbox}[
  enhanced,
  breakable,
  colback=white,
  colframe=FinanceBlue,
  boxrule=0.7pt,
  sharp corners,
  borderline west={2.2pt}{0pt}{FinanceBlue},
  title=\textbf{财务展望：现金流权益与投票权不必一致},
  coltitle=white,
  colbacktitle=FinanceBlue,
  fonttitle=\bfseries,
  before skip=1.2em,
  after skip=1.2em
]

某些类别的股份可能仅占公司总现金流的很小一部分，但它们却拥有较大的表决权。现金流权利决定了参与剩余经济价值分配的权利，而表决权则决定了影响公司决策的能力。

\end{tcolorbox}

\subsection{公司行动与所有权请求权的变化}

公司行为可以改变企业持有的现金数量、存续股票的数量，或者改变股东之间的权利分配。这些行为的直接影响可以通过分析现金是进入还是离开企业、股票是增加还是减少，以及所有权份额或权利是否发生变化来分析。

新的股权发行会增加存续的股数，并可能改变现有的股权份额。假设一家公司有\(N_0\)存续股份和发行股份\(N_1\)新股。持有这些股票的投资者……\(n\)在发行前持有的股份只占一小部分。$\alpha_0=n/N_0$. 如果投资者不参与本次发行，那么其持有股份比例将变为……$\alpha_0=n/(N_0+N_1)$  

投资者的持股比例受到稀释，但比例稀释并不必然意味着经济价值损失。如果新股按公平价格发行，公司收到的现金在发行时使其价值大致增加相应筹资金额，因此较小的持股比例可能对应更大的总股权价值。如果以异常低价向特定投资者发行，价值则可能从未参与发行的股东转移给新投资者。

在LaTeX命令、公式、交叉引用以及数字方面，保持原有的格式即可。  

股票发行允许现有股东购买新股份，如果他们参与这一计划，就可以保持其持股比例。而不参与该计划的股东则可能会面临股份比例减少的情况；不过，可转让权利可以为这些股东提供一定的补偿，以弥补因放弃认购特权而造成的损失。

\begin{example}{百分比稀释与数值转移}{new-equity-dilution}

该公司有8000万股存续股票，每股交易价格为10元人民币左右，因此其市场股权价值约为8亿元人民币。一位投资者持有800万股股票，占该公司总股本的10\%。

该公司发行了2000万股新股票，而该投资者并未参与此次发行。因此，该投资者的持股比例下降到了……

\[
\frac{8}{100}=8\%.
\]

如果新发行的股票价格为人民币10元，那么公司筹集了约2亿元人民币的资金。假设其他价值没有变化，发行后的股权价值将变为10亿元人民币，而股价则保持在10元人民币。投资者的持股比例会下降，但800万股股票的价值仍然为8000万元人民币。

如果新发行的股份以人民币约5元的价格分配给选定的投资者，那么公司仅筹集了约1亿元人民币的资金。在相同的假设下，发行后的股权价值约为9亿元人民币，这意味着每股的股价为

\[
\frac{900}{100}=\text{RMB }9.
\]

该投资者的股份价值约为7200万元人民币。在这种情况下，这种低于市价的发行方式实际上是将价值从现有股东转移到了新投资者手中。

因此，百分比稀释以及经济价值的损失是两种不同的情况。经济效应取决于所有权比例的变动以及发行价格的因素。

\end{example}

股利和股票回购都是从公司向股东分配现金，但它们对所有权请求的影响不同。现金股利是将资源从公司转移到股东手中，而不会改变存续股份的数量。如果一家公司每股支付股利\(D\)，并且没有其他相关信息，那么在股票开始交易时，没有该股利的股东们的股票价格应该会下降。忽略税收和市场摩擦因素的话……

\begin{equation}
P_{\mathrm{ex}}
\approx
P_{\mathrm{cum}}-D,
\label{eq:ch9-ex-dividend-price}
\end{equation}

其中，\(P_{\mathrm{cum}}\)表示包含股利权利的股票价格，而\(P_{\mathrm{ex}}\)则代表除股利权利后的股票价格。股东会获得现金，同时公司的现金持有量也会减少。

股份回购也将公司现金转移给股东，但会减少已发行在外股份。出售股份的股东获得现金，在其他条件不变时，保留股份的股东持有剩余股权的比例提高。回购可能通过支付价格、交易传递的信息、税收或资本结构变化影响价值。仅减少股份数量，并不会创造经济价值。

股票分割改变了股份的数量，但并未转移现金或改变所有权比例。在 \(m\) 对 \(n\) 的分割中，股份数量会发生变化。

\begin{equation}
N'=\frac{m}{n}N,
\label{eq:ch9-stock-split-shares}
\end{equation}
如果没有其他因素的影响，那么每股的价格大约会调整为

\begin{equation}
P'=\frac{n}{m}\times P.
\label{eq:ch9-stock-split-price}
\end{equation}

因此，如果拆分操作并未传递任何新的信息，也不会对市值产生任何影响，那么总市值将保持不变。在输出结果时，请保留所有LaTeX命令、公式表达、交叉引用以及数字等内容。

\begin{table}[H]
\centering

\caption{选定的企业行为及其直接的影响}
\label{tab:ch9-corporate-actions}

\begin{tabularx}{\textwidth}{@{}p{2.8cm}Y Y@{}}

\toprule

\textbf{操作}
& \textbf{立即更改}
& \textbf{关键的经济问题} \\

\midrule

现金股利
& 现金分配给股东；存续股份保持不变。在除息日，股价通常会下降。
& 这种分配方式如何影响公司保留的资源和股东价值？\\

股份回购  
& 现金流出公司，导致存续股份普遍减少。  
& 以什么价格回购股份？这些操作如何影响股份持有比例和股东价值？\\

股票分割
& 存续股份增加，每股价格相应调整；总股权价值和持股比例保持不变，其他条件相同。
& 这种分割是否传递了信息或影响了交易条件？\\

新股权发行  
& 现金进入公司，存续股份增加，可能会稀释现有股东的持股比例。  
& 新股份以什么价格发行？这种发行对现有股东的经济价值有何影响？\\

\bottomrule

\end{tabularx}

\end{table}

\section{股票市场架构与价格发现}

股票市场的运作具有双重功能。一级市场允许企业筹集股权资本，而二级市场则使得现有股份能够在投资者之间转让。在次级交易中，基于投资者的信息和预期的竞争性买卖订单决定了市场价格，这一价格由交易所、经纪商、清算系统和存托机构共同支撑。

\subsection{股权发行与一级市场}

一家公司通过向投资者发行股票在一级市场上筹集新的股权资本。首次公开募股是指首次向公众投资者发行股票，而成熟股权融资则是指在上市之后继续筹集额外的公众股权。权益发行和私人配售则是发行新股票的另一种方式。

股权发行需要确定新请求权的条款和价格。承销商决定提供多少股份以及这些股份附带哪些权利，而投资者则评估公司的预期现金流、风险、治理情况以及资金的使用方式。承销商、簿记人以及财务顾问等中介机构可以协助进行尽职调查、信息披露、簿记建档以及股份发行工作。他们的角色是促进发行过程，并将新股份与投资者的需求相匹配。

公开发行可能同时包含公司新发行股份和现有股东出售的股份。前者是公司新发行的股份，向公司注入新资金；后者由现有股东出售，只转移所有权，出售所得不进入公司。解释发行筹资额时，这一区别十分重要。

\begin{figure}[H]
\centering

\begin{tikzpicture}[
  >=Latex,
  box/.style={
    draw,
    rounded corners=2pt,
    line width=0.7pt,
    align=center,
    text width=1.9cm,
    minimum height=0.9cm,
    inner sep=5pt
  },
  cash/.style={
    -Latex,
    line width=0.9pt,
    FinanceBlue
  },
  claim/.style={
    -Latex,
    dashed,
    line width=0.9pt,
    FinanceBlue
  },
  lab/.style={
    font=\small,
    align=center,
    fill=white,
    inner sep=1.5pt
  }
]

% Participants
\node[box] (firm) at (-1.0,0)
  {公司};

\node[box] (market) at (4.5,0)
  {一级市场};

\node[box] (investors) at (10.0,0)
  {投资者};

% New equity claims: company to investors
\draw[claim]
  ([yshift=5pt]firm.east) --
  node[lab, above=3pt] {新的股权请求权}
  ([yshift=5pt]market.west);

\draw[claim]
  ([yshift=5pt]market.east) --
  node[lab, above=3pt] {新的股权请求权}
  ([yshift=5pt]investors.west);

% Capital: investors to company
\draw[cash]
  ([yshift=-5pt]investors.west) --
  node[lab, below=3pt] {大写字母}
  ([yshift=-5pt]market.east);

\draw[cash]
  ([yshift=-5pt]market.west) --
  node[lab, below=3pt] {大写字母}
  ([yshift=-5pt]firm.east);

\end{tikzpicture}

\caption{一级市场的资本与股权请求权流动}
\label{fig:ch9-primary-equity}
\end{figure}

发行价格会影响所筹集的资本数量，以及现有股东与新股东之间的所有权分配。如果发行价格相对于投资者的估值来说过高，那么本次发行可能会吸引不足的需求。而如果发行价格过低，发行人可能筹集到的资金会少于原本能够获得的资金。初始定价折扣可能反映了估值的不确定性、信息不对称以及发行风险。

\subsection{交易所与其他交易场所}

交易所为证券的上市和交易提供规则与基础设施。上市与交易虽然相关，但却是两个不同的过程。一家公司在某个交易所上市后，其股票的交易也可以通过其他允许的交易平台进行。

在共同的交易规则、准入安排和市场数据系统的框架下，交易商将订单汇集在一起进行交易。现代股权市场还可能包含其他交易系统，例如交叉网络和暗池，以及以交易商为基础的交易平台。

假设有一家公司的股票在纽约证券交易所上市。投资者并不一定要在纽约证券交易所进行这些股票的交易。根据相关的规则和订单处理机制，经纪人可以将这些订单提交给其他交易所或采用替代交易系统来处理，比如暗池交易系统。因此，同样的股票可以在多个交易平台上进行交易，每个平台都有自己独特的订单流和流动性状况。

多个交易场所会增加执行服务的竞争，同时允许不同的交易机制共存。不过，这也可能导致订单流出现混乱，使得市场中的流动性难以被有效监控。因此，交易场所的结构决定了订单之间的相互作用方式、流动性的集中程度以及价格的形成方式。

交易场所在更广泛的市场基础设施中运作。经纪商负责传递订单，清算机构负责建立结算机制，并在适用的情况下处理与交易相关的义务。而证券保管机构则负责维护所有权记录，结算系统则负责完成现金和证券的转移。电子交易改变了这些过程的速度和组织方式，但并未改变其基本经济功能。

\subsection{交易机制、订单与市场流动性}

一种交易机制决定了买卖订单如何相互作用。在订单驱动的市场中，订单被提交到限价单簿中，当兼容的订单相遇时，交易就会发生。而在交易商市场里，交易商会报出他们愿意买入和卖出的价格。许多现代股权市场结合了这两种结构的元素。

市价单优先考虑按可得最优价格成交，但受可用数量限制。限价单规定买方愿意支付的最高价格或卖方愿意接受的最低价格，提供更强的价格控制，却可能无法成交。市价单通常需求现有流动性，而不能立即成交的限价单可为其他交易者供给流动性。

如果最高出价和最低出价都是\(P_{\mathrm{bid}}\)以及那些使用LaTeX命令、公式表达、交叉引用和数字的格式。\(P_{\mathrm{ask}}\)，所引用的买卖价差为

\begin{equation*}
s=P_{\mathrm{ask}}-P_{\mathrm{bid}}.
\label{eq:ch9-bid-ask-spread}
\end{equation*}

中点是

\begin{equation*}
M=\frac{P_{\mathrm{ask}}+P_{\mathrm{bid}}}{2},
\label{eq:ch9-midquote}
\end{equation*}
以及相应的比例区间范围如下：

\begin{equation*}
s_{\mathrm{rel}}
=
\frac{P_{\mathrm{ask}}-P_{\mathrm{bid}}}{M}.
\label{eq:ch9-relative-spread}
\end{equation*}

价差是衡量流动性的一方面。深度指的是当前价格附近可获得的交易量，而价格影响则反映了执行价格对订单大小的响应程度。对于小额交易来说，市场可能会呈现较小的价差，但对于大额订单而言，市场的深度则可能有限。

市场流动性会随着市场状况的变化而波动。当不确定性或逆向选择风险增加时，流动性提供者可能会报出更宽的价差，展示更少的深度，或者撤回订单。因此，与更了解市场的投资者进行交易的可能性，使得信息条件与市场流动性紧密相连。

交易限制也会影响信息和投资者观点如何进入订单流。卖空机制允许投资者借入股票进行出售，之后再回购这些股票以偿还借款。这样一来，即使投资者并不实际持有这些股票，他们仍然可以通过卖出订单来表达负面信息或预期。不过，相关的限制或较高的借款成本可能会限制这种交易方式的效果。因此，卖空行为确实会对订单流、流动性以及价格形成产生影响，但其具体影响程度则取决于市场状况和交易规则。

\subsection{信息、价格与价格发现}

股票价格是股权请求权的实际交易价格。当投资者对未来现金流或必要报酬率的预期发生变动，或者与请求权相关的权利发生变化时，股票价格就会发生变化。

用 \(\mathcal{I}_t\) 表示在时间 \(t\) 时可获得的信息。一种简化的基于股利的估值可以表示为

\begin{equation}
P_t
=
\sum_{j=1}^{\infty}
\frac{\mathrm{E}[D_{t+j}\mid\mathcal{I}_t]}
{(1+r_{e,t})^j},
\label{eq:ch9-information-price}
\end{equation}
其中，\(r_{e,t}\)被视为从时间\(t\)开始所需的报酬率。新的信息可能会改变预期的未来股利、所需报酬率，或者两者都会发生变化，因此也会改变投资者对这支股票的估值。

价格形成是指信息和投资者评估通过交易被反映到市场价格过程中的过程。不同投资者的信息、模型、预期以及流动性需求各不相同，因此他们的订单只能反映出他们部分已知的信息。因此，价格形成虽然能够提供有用信息，但并不一定是即时或无需成本的。图~\ref{fig:ch9-price-discovery}展示了信息及投资者评估如何被反映到观察到的市场价格中，同时也说明了市场摩擦对这一过程的影响。


\begin{figure}[H]

\centering

\begin{tikzpicture}[
  >=Latex,
  box/.style={
    draw,
    rounded corners,
    align=center,
    text width=2.55cm,
    minimum height=0.9cm
  },
  flow/.style={
    -Latex,
    line width=0.85pt,
    FinanceBlue
  }
]

\node[box] (info) at (0,1.4)
  {公称和私称\\信息};

\node[box] (beliefs) at (3.3,1.4)
  {现金流与风险评估\\};

\node[box] (orders) at (6.6,1.4)
  {命令、公式、交叉引用和数字部分保持不变。\\交易};

\node[box] (price) at (9.9,1.4)
  {已观察到\\市场价格};

\draw[flow] (info) -- (beliefs);
\draw[flow] (beliefs) -- (orders);
\draw[flow] (orders) -- (price);

\node[
  draw,
  rounded corners,
  align=center,
  minimum width=5.8cm,
  minimum height=0.75cm
] (frictions) at (6.6,-0.3)
  {流动性、资金供应以及市场中的摩擦因素。};

\draw[flow] (frictions.north) -- (orders.south);

\end{tikzpicture}

\caption{信息、交易与价格发现}
\label{fig:ch9-price-discovery}
\end{figure}

有效市场假说为评估市场价格如何吸收信息提供基准，其含义取决于所考虑的信息集。在该假说下，扣除交易和信息成本后，投资者不应能够系统地利用相应市场有效形式所涵盖的信息，获得异常的风险调整后回报。

根据所考虑的信息类型，假设通常可以用三种方式来表达。弱形式效率关注的是过去的价格和交易活动中包含的信息。半强形式效率则扩展了信息的范围，包括了公开可用的信息，比如财务报告和企业公告等。而强形式效率则同时考虑了公开信息和私有信息，这一标准更可以被视为一种要求较高的基准，而非大多数市场的真实描述。

价格发现仍然受到经济摩擦的影响。获取信息可能成本高昂，投资者之间可能存在分歧，此外，流动性状况、资金限制以及套利的限制也会对观察到的价格产生影响。如果价格完全反映了所有难以获取的信息，那么投资者就没有动力去承担获取这些信息的成本。

\begin{tcolorbox}[
  enhanced,
  breakable,
  colback=white,
  colframe=FinanceBlue,
  boxrule=0.7pt,
  sharp corners,
  borderline west={2.2pt}{0pt}{FinanceBlue},
  title=\textbf{财务展望：价格发现并非完美估值},
  coltitle=white,
  colbacktitle=FinanceBlue,
  fonttitle=\bfseries,
  before skip=1.2em,
  after skip=1.2em
]

观察到的市场价格反映了投资者在当前市场条件下的信息、预期以及评估。这些价格虽然能反映经济价值，但它们并非对客观存在的基本价值的直接体现。

\end{tcolorbox}

\section{股权价值与风险}

股权价值反映了可供股东使用的现金流的现值，同时考虑了这些现金流的风险因素。股权估值有三个核心要素：预期股东现金流、再投资对未来增长的影响，以及与股权风险相一致的必要报酬率。

\subsection{股权现金流与估值输入}

应当区分账面股权和市场股权。账面股权是一种基于已确认的资产和负债得出的会计计量方式，而市场股权则是存续的请求权的总市值。这两者并不一定要一致，因为市值反映了对未来现金流、增长机会、风险以及与股份相关的权利的预期。

股利是指公司直接支付给股东的现金。股票回购则是一种另一种形式的支付方式：公司向出售股票的股东支付现金，同时使存续的股数减少。因此，股利和股票回购以不同的方式分配价值。

留存收益在保留时并不作为对股东的分配。其对股权价值的贡献取决于公司内部保留的资源所产生的未来现金流。只有当再投资能够支持具有足够价值的未来现金流时，保留收益才会创造价值。

对于计划持有股票直到日期\(T\)的投资者来说，当前的价值可以表示为在持有期间预期获得的股利与预期出售价格的现值之和。

\begin{equation}
P_0
=
\sum_{t=1}^{T}
\frac{\mathrm{E}[D_t]}{(1+r_e)^t}
+
\frac{\mathrm{E}[P_T]}{(1+r_e)^T},
\label{eq:ch9-finite-horizon-equity-value}
\end{equation}
其中，\(r_e\) 是在恒定利率简化条件下，股权请求权的必要报酬率。

如果终端价格被视为在日期$T$之后预期股利的现值，那么无限延长估值范围，就可以得到股利折现表示法：

\begin{equation}
P_0
=
\sum_{t=1}^{\infty}
\frac{\mathrm{E}[D_t]}{(1+r_e)^t}.
\label{eq:ch9-dividend-discount-model}
\end{equation}


对于股权估值，分析师可以预测自由现金流至股权，即\(FCFE_t\)，这是扣除运营需求、税收、再投资以及净债务融资后，可供股权持有人使用的现金流。一个简化的估值表达式为

\begin{equation}
V_E
=
\sum_{t=1}^{\infty}
\frac{\mathrm{E}[FCFE_t]}{(1+r_e)^t}.
\label{eq:ch9-fcfe-valuation}
\end{equation}

现金流指标与贴现率应对应同一请求权。可供全部资本提供者获得的现金流，需要采用适合整个企业的贴现率；仅可供股权持有人获得的现金流，则需要采用股权必要报酬率。

可比公司的市场价格提供了额外的估值基准。相对估值则使用诸如市盈率、市净率等指标来比较类似的请求权。当考虑到增长、盈利能力、杠杆率、风险以及会计计量方面的差异时，这种比较更具参考价值。

\subsection{增长预期与市场价值}

增长通过未来现金流及实现增长所需的再投资影响股权价值。预期销售或利润提高，本身并不意味着价值提高。新增股权投入所获得的回报超过股东要求的回报时，再投资才创造价值。

股利支付比率将预期收益与下一期的股利联系起来。假设下一期的每股收益为$E_1$，公司保留一部分收益$k$，并将剩余部分作为股利分配。那么

\begin{equation}
D_1=(1-k)E_1.
\label{eq:ch9-dividend-payout}
\end{equation}

如果增长是通过留存收益来融资的，那么就不会发行新的普通股股权，而期初的股权回报率保持稳定。在这种情况下，一个简化的可持续增长关系为

\begin{equation}
g=k\,ROE.
\label{eq:ch9-sustainable-growth}
\end{equation}

如果股利以恒定速率
$g$ 无限期地增长，那么戈登增长模型，也被称为恒定增长股利折现模型，会表示当前的股价为

\begin{equation}
P_0=\frac{D_1}{r_e-g},
\qquad
r_e>g.
\label{eq:ch9-gordon-growth}
\end{equation}

代入收益和增长关系后，结果如下：

\begin{equation}
P_0
=
\frac{(1-k)E_1}
{r_e-k\,ROE}.
\label{eq:ch9-growth-reinvestment-value}
\end{equation}

较高的留存比率虽然会减少当前的支付额，但能够提升未来的增长潜力。其对价值的影响取决于所保留资源所获得的报酬率与股权必要报酬率之间的关系。

\begin{example}{增长与价值创造}{growth-reinvestment-value}

假设下一期的每股收益为10，留存率为40\%，而股权的必要报酬率为10\%。

如果留存收益能够保持稳定的水平\(ROE\)15中的…\%,

\[
g=0.40\times 0.15=6\%,
\]
与
\[
P_0=\frac{(1-0.40)\times 10}{0.10-0.06}=\frac{6}{0.04}=150.
\]

如果股权回报为8\%，

\[
g=0.40\times 0.08=3.2\%,
\]
与
\[
P_0=\frac{(1-0.40)\times 10}{0.10-0.032}=\frac{6}{0.068}\approx88.24.
\]

在这种简化的设定中，当再投资获得的收益超过必要报酬率时，增长会增加价值；而当再投资的收益低于必要报酬率时，增长则会降低价值。

\end{example}

这一原则不仅适用于恒定增长模型。快速增长可能需要大量投资，而适度的当前收益在保留资源能够产生足够强劲的未来现金流时，仍然可以体现出较高的价值。增长的预期也会影响价值的实现时间。当预期中的更多现金流发生在更远的未来时，股权价值对必要报酬率的变化就变得更加敏感。

\begin{tcolorbox}[
  enhanced,
  breakable,
  colback=white,
  colframe=FinanceBlue,
  boxrule=0.7pt,
  sharp corners,
  borderline west={2.2pt}{0pt}{FinanceBlue},
  title=\textbf{财务展望：增长并不等同于价值创造},
  coltitle=white,
  colbacktitle=FinanceBlue,
  fonttitle=\bfseries,
  before skip=1.2em,
  after skip=1.2em
]

只有当所需再投资产生的回报足够高时，增长才能提升股权价值，而这一回报必须高于股东所要求的回报。

\end{tcolorbox}

\subsection{股权风险溢价与必要报酬率}

投资者在某一时期内持有的股票的实际回报，是所收到的股利与市场价格变化的总和。如果初始价格为 \(P_0\)，期末价格为 \(P_1\)，股利为 \(D_1\)，那么实际持有的期间回报就是

\begin{equation}
r
=
\frac{P_1-P_0+D_1}{P_0}.
\label{eq:ch9-equity-realized-return}
\end{equation}

必要报酬率不同于已实现回报率。它是投资者基于股权请求权风险及其他可得投资回报，要求为持有该请求权获得的回报。能够通过多项投资分散的风险，通常比与广泛经济条件相关的风险需要更少补偿。

令 \(r_M\) 表示广泛市场投资组合的回报，而 \(r_f\) 表示无风险利率。市场股权风险溢价为

\begin{equation}
ERP
=
\mathrm{E}[r_M]-r_f.
\label{eq:ch9-equity-risk-premium}
\end{equation}

股权风险溢价具有前瞻性。历史股票已实现回报超过无风险回报的平均值，可为过去结果提供证据，却未必等于特定日期投资者要求的溢价。

资本资产定价模型是一种被广泛使用的基准方法，用于将股票的必要报酬率与系统性风险联系起来。股票\(i\)的贝塔值为

\begin{equation}
\beta_i
=
\frac{\mathrm{Cov}(r_i,r_M)}
{\mathrm{Var}(r_M)}.
\label{eq:ch9-beta}
\end{equation}

在资本资产定价模型假设下，必要报酬率为

\begin{equation}
r_i
=
r_f
+
\beta_i\left(\mathrm{E}[r_M]-r_f\right).
\label{eq:ch9-capm}
\end{equation}

贝塔系数衡量了股票对市场整体回报变动的敏感度。贝塔系数越高，意味着在资本资产定价模型中的风险溢价要求也就越高。

\begin{example}{A CAPM 必要报酬率}{capm-required-return}

假设无风险利率为3\%，预期市场股权风险溢价为6\%，某股票的贝塔值为1.2。那么CAPM的必要报酬率为

\[
r_i
=
3\%+1.2(6\%)
=
10.2\%.
\]

该估计结果是基于CAPM模型以及所使用的输入数据而得出的。

\end{example}

CAPM~是一个基准，而非对预期回报的完整描述。其更广泛的经济含义是：与全市场不利条件相关的风险，在保护最有价值时难以分散，因而通常需要更高补偿。贝塔通过与市场回报的协方差表示这种风险暴露，其他资产定价模型则允许存在额外的定价风险来源。

%========================================================
% Summary
%========================================================

\TailModule{summary}{本章小结}

普通股是一种剩余所有权请求权。股东在优先 contractual 请求权得到满足之后获得价值，并参与剩余的增值收益；而有限责任通常可以防止完全支付的股份所带来的损失超出该股份所投入的价值。股权融资可以提供资本，而不需要产生固定的 contractual 利息或本金支付，并且新投资者可以获得剩余现金流的权利，通常还有表决权。即使新发行的股份会稀释股东的百分比所有权，但它也会增加公司可用的资源。

公司治理分配对企业资源的控制权。股东通常通过投票和选举董事来间接行使控制权，而董事会则负责任命和监督管理层。分散化的所有权可能会削弱监督的激励作用，而集中的所有权则可以加强监督，同时增加控股股东与少数股东之间的潜在冲突。不同的股票类别可以区分现金流权和表决权。公司的各种行为，如股利分配、股票回购、股票分割和新股发行，都会以不同的方式影响现金、股票数量以及各股东的请求权。

股票市场的运作既具有融资功能，也具备交易功能。首次发行可以为公司引入新的资本，而二次交易则使得现有股份在投资者之间流转。交易所及其他交易平台为订单的传递提供了机制，而市场的流动性则取决于价差、深度、价格影响、信息状况，以及市场参与者提供流动性的意愿。价格的形成反映了市场上的信息、预期、风险评估以及流动性状况，但这并不意味着每一个观察到的价格都等同于某种客观认定的基本价值。

股权估值要求股东现金流与用于折现的必要报酬率之间保持一致性。股利和回购会释放现金，而留存收益则通过未来再投资的现金流为股权创造价值。当预期再投资的报酬率远高于股权必要报酬率时，股权价值会上升。市场股权风险溢价代表了广泛股权投资的预期超额回报，而CAPM则提供了一个基准，将个股的必要报酬率与其系统性风险联系起来。

%========================================================
% Key Terms
%========================================================

\TailModule{keyterms}{关键术语}

普通股权（Common equity）；剩余请求权（residual claim）；有限责任（limited liability）；持股比例稀释（ownership dilution）；投票权（voting rights）；董事会（board of directors）；公司治理（corporate governance）；代理问题（agency problem）；控股股东（controlling shareholder）；中小股东（minority shareholder）；股份类别（share class）；优先股（preferred share）；不同投票权（weighted voting rights）；现金股利（cash dividend）；股份回购（share repurchase）；股票拆分（stock split）；配股（rights issue）；首次公开发行（initial public offering）；上市后增发（seasoned equity offering）；公司新发行股份（primary shares）；现有股东出售股份（secondary shares）；交易所（exchange）；其他交易场所（alternative trading venue）；市价单（market order）；限价单（limit order）；买卖价差（bid--ask spread）；市场流动性（market liquidity）；价格发现（price discovery）；信息有效性（informational efficiency）；股利贴现模型（dividend discount model）；股权自由现金流（free cash flow to equity）；留存比例（retention ratio）；可持续增长（sustainable growth）；股权风险溢价（equity risk premium）；贝塔（beta）；资本资产定价模型（Capital asset pricing model）。

%========================================================
% Financial AI Lab
%========================================================

\TailModule{aiprompt}{金融~AI~实验室}

\textbf{实验：对人工智能的股权估值和治理情况进行审计与比较}

选择两家属于同一大行业的上市公司。让人工智能系统使用简单的价格倍数法，以及基于自由现金流与股权或股利的评估方法来比较它们的股权估值。要求该系统明确指出所评估的股类、相关的存续股份数量、预期的增长假设、必要报酬率，以及每个事实输入的来源。

将结果分为四个阶段进行审核。

\begin{enumerate}

\item 验证市值和每股估值所使用的股份数量是否正确，以及每个估值倍数是否与相应的股权指标相匹配。

\item 研究估值对增长假设和必要报酬率的响应。至少测试一种不同的增长率和一种不同的必要报酬率，或者，在适用的情况下，测试股权风险溢价。

\item 请比较所选股票类别所附带의投票权和现金流权利。确定是采用“一股一票”的原则，还是通过治理结构来分离投票权与经济所有权的关系。

\item 请对照公司的公开信息或官方的消息来源来核实每一个事实性声明。至少找出一个声明，人工智能提供了合理的解释，但缺乏足够的支持证据。

\end{enumerate}

%========================================================
% Discussion Questions
%========================================================

\TailModule{questions}{讨论题}

\begin{enumerate}

\item 为什么有限责任并不能使普通股权成为低风险的请求权呢？需要区分的是，从拥有股票本身所产生的最大损失，以及股票市场价值的变化。

\item 如何使得集中所有权能够减少某个代理问题，同时增加另一个代理问题的解决呢？

\item 为什么在利好消息公布后，即使公司没有从股市获得新的资金注入，股价仍然能够大幅上涨呢？

\item 为什么更快的盈利增长并不一定意味着更高的股权价值？请结合再投资和股权必要报酬率来阐述你的答案。

\item 为什么将历史平均实现的市场回报与无风险回报之间的差额自动视为当前预期的股权风险溢价是不合理的呢？

\end{enumerate}

%========================================================
% Applied Problems
%========================================================

\TailModule{appprobs}{应用题}

\begin{enumerate}

\item \textbf{剩余支付与有限责任。}

一家以债务和股权为融资方式的公司，在下一个时期将产生160或70的现金流，这两种情况出现的概率都是相等的。承诺的债务偿还金额为100。

\begin{enumerate}

\item 计算每个状态下的债务和股权支付情况。

\item 计算股权部分的预期支付。

\item 解释股权支付结构中的哪一部分体现了有限责任。

\item 解释为什么在发行债务之后，股东可能会从公司现金流风险性的转变中获益，而债权人则可能因此受损。

\end{enumerate}

\item \textbf{投票权和现金流控制权。}

有一家公司拥有2000万股A类股票，每股对应的投票数为10票；同时，该公司还有1.8亿股B类股票，每股对应的投票数为一票。假设这两类股票的每股现金流权利是相同的。

\begin{enumerate}

\item 计算该公司总现金流中A类资产的份额。

\item 计算A类股份在总投票权中的份额。

\item 请解释为什么这两个百分比之间存在差异。

\item 明确指出这种结构的一个潜在优势以及一个潜在的治理风险。

\end{enumerate}

\item \textbf{市价订单和限价订单。}

该股票的最佳买入报价为24.96美元，对应的股票数量为6,000股；而最佳卖出报价则为25.04美元，对应的股票数量则是4,000股。

\begin{enumerate}

\item 计算报价差值和中间值。

\item 计算相应的报价价差比例。

\item 如果报价仍然有效，那么以什么价格可以执行一笔购买2000股的市场订单呢？

\item 如果下达了一个买入限价订单，订单数量是2000股，价格目标是25.00元，而没有任何卖方愿意在25.00元或更低的价格出售股票，那么这个订单在限价订单簿中将会如何处理呢？

\end{enumerate}

\item \textbf{增长与再投资。}

一家公司预计下一期的每股收益为8，其股权必要报酬率为11\%。该公司将保留30\%的收益。假设增长是通过留存收益来融资的，没有新的普通股股本被发行，且\(ROE\)保持在16\%的水平不变。

\begin{enumerate}

\item 计算可持续增长率。

\item 计算下一周期的每股股利。

\item 使用恒定增长模型来计算股票的价值。

\item 如果需要进行重新计算，请点击此处。保留LaTeX命令、公式、交叉引用和数字。\(ROE\)变为9\%虽然保留率保持不变。

\item 请说明估值变化的方向。

\end{enumerate}

\item \textbf{股权风险与必要报酬率。}

无风险利率为2.5\%，预期市场股权风险溢价为5.5\%。

\begin{enumerate}

\item 计算贝塔值为0.8的股票的CAPM必要报酬率。

\item 计算贝塔值为1.4的股票的必要报酬率。

\item 解释在资本资产定价模型（CAPM）中，为什么高贝塔值的股票具有更高的必要报酬率。

\item 请给出一个理由，说明为什么分析师可能不愿意仅依赖CAPM估值来进行决策。

\end{enumerate}

\end{enumerate}

%========================================================
% Data Application
%========================================================

\TailModule{dataapp}{数据应用}

\textbf{中国股权市场的信息与价格发现}

请选择一家在上海证券交易所或深圳证券交易所上市的中国A股公司，并找到一份包含经济相关新信息的企业公告，例如收益公告、利润预警、重大投资决策、股利提案、股票回购计划或新的股权发行。以交易所披露平台和该公司的正式公告作为获取这些信息的首要来源。

收集在披露事件前后几天的股票每日收盘价，同时记录同一时期的市场指数。计算每天的股票回报率以及简单的市场调整后的异常回报值。

\begin{equation*}
r_t=r_{i,t}-r_{M,t},
\end{equation*}

其中，\(r_{i,t}\) 表示第 \(t\) 天的股票回报，而 \(r_{M,t}\) 则表示市场指数的回报。

确定信息首次公开披露的具体时间。如果披露信息是在市场关闭后发布的，那么下一个交易日应被视为投资者可以通过交易进行响应的第一天。绘制股票价格的走势图，并报告在披露信息发布后的市场调整收益情况。判断最大的价格波动发生在披露信息后的第一个交易机会之前、期间还是之后。

谨慎解读这一结果。价格反应可能反映了对预期现金流或必要报酬率的调整，或者两者都有。同时发布的全市场新闻以及其他公司特定的信息也可能影响观察到的回报。较短的事件窗口和每日收盘价本身并不能证明因果关系或日内价格发现的速度。

请列出该公司的名称、股票代码、披露标题、披露日期和时间、样本周期、基准指数以及数据来源。同时说明在得出市场对这些信息反应迟缓或过度反应的结论之前，还需要哪些额外的证据。

%========================================================
% Case Analysis
%========================================================

\TailModule{caseanalysis}{案例分析}

\textbf{案例：小米与香港的加权投票权}

小米公司在2018年7月9日开始在香港上市，成为在2018年上市制度改革后首批采用加权投票权结构在香港上市的公司。小米的A类股票在大多数股东决议中拥有每股十票的权利，而B类股票则拥有每股一票的权利。对于某些特定保留事项，每股股票拥有一票表决权。目前，小米公开上市的股票属于B类股票。

这种结构将投票权与经济所有权分离开来。它允许创始人在从公众投资者那里筹集到股权后，仍然对公司的决策拥有相当大的影响力。香港的上市规则允许采用特定的加权投票结构，但前提是必须满足相关资格要求、持续履行义务，并为非加权投票的股东提供相应的保护措施。

从融资和治理的角度分析该结构。为什么创始人可能重视能够筹集到公开股权，而不放弃相应的投票控制权？集中的投票权力如何支持战略上的连续性或长期投资？当投票控制权超过现金流所有权时，可能会出现哪些代理问题？当外部股东的投票权力较弱时，哪些投资者保护机制变得更为重要？

那么考虑一下市场价值的影响。解释为什么小米公开发行B类股票的市场价格不仅反映了预期的剩余现金流，还反映了与投资相关的治理权和控制权。如果预期运营现金流最初没有变化，那么表决权的变化是否一定会使市场价值保持不变？解释控制权的变化、代理成本、未来现金流的分配以及必要报酬率如何影响股权价值。

\emph{清算：香港交易及结算所有限公司，2018年7月4日发布的《HKEX将为投资者提供更多选择，助力小米股票上市》；小米公司，2018年招股说明书及2018年中期报告；香港证券交易所上市规则，关于加权投票权的第8A章。}

%========================================================
% Further Reading
%========================================================

\TailModule{furtherread}{延伸阅读}

\begin{itemize}

\item Gordon, Myron J. (1959),
``Dividends, Earnings, and Stock Prices（股利、收益与股票价格）,''
\emph{The Review of Economics and Statistics}, 41(2, Part 1), 99--105.

\item Fama, Eugene F. (1970),
``Efficient Capital Markets: A Review of Theory and Empirical Work（有效资本市场：理论与实证研究综述）,''
\emph{The Journal of Finance}, 25(2), 383--417.

\item Grossman, Sanford J., and Joseph E. Stiglitz (1980),
``On the Impossibility of Informationally Efficient Markets（论信息有效市场的不可能性）,''
\emph{American Economic Review}, 70(3), 393--408.

\item Kyle, Albert S. (1985),
``Continuous Auctions and Insider Trading（连续拍卖与内幕交易）,''
\emph{Econometrica}, 53(6), 1315--1335.

\item La Porta, Rafael, Florencio Lopez-de-Silanes, Andrei Shleifer, and Robert W. Vishny (2000),
``Investor Protection and Corporate Governance（投资者保护与公司治理）,''
\emph{Journal of Financial Economics}, 58(1--2), 3--27.

\end{itemize}

\ChapterTail{summary,keyterms,aiprompt,questions,appprobs,dataapp,caseanalysis,furtherread}
